\documentclass[11pt]{article}
\usepackage[a4paper,margin=1in]{geometry}
\usepackage{verbatim}
\usepackage{parskip}
\usepackage{hyperref}
\usepackage{enumitem}
\usepackage{tabularx}

\title{DTree Public API Reference}
\author{}
\date{}

\newcolumntype{L}{>{\raggedright\arraybackslash}X}
\newcommand{\apimethod}[3]{%
	\noindent\begin{tabularx}{\textwidth}{@{}lL@{}}
		\textbf{#1 -- } & #3 \\
		& \texttt{#2} \\
	\end{tabularx}
	\vspace{-0.3\baselineskip}
	\begin{itemize}[leftmargin=2em, label=--]%
	}
	
	\begin{document}
		\maketitle
		
		\section*{Overview}
		
		The \texttt{DTree} class represents a symbolic directed tree with a unique root and internal support for dual projections (Q1 and Q2). Each node is identified by a unique string. Edges are added incrementally, and projections are updated automatically. Structural operations preserve the tree’s topological validity.
		
		\section*{Constructor}
		
		\apimethod{Constructor}{Create a new tree with a given identifier}{DTree(string id)}
		\item Initializes a new directed tree with the given identifier.
		\item The root is determined implicitly as the source of the first inserted edge.
		\item Q1 and Q2 projections are initialized automatically.
	\end{itemize}
	
	\section*{Core Methods}
	
	\apimethod{Insert}{Insert a new edge into the tree}{void Insert(string parent, string child)}
	\item Adds a directed edge from \texttt{parent} to \texttt{child}.
	\item The \texttt{parent} must already exist in the tree.
	\item The \texttt{child} must not already exist.
	\item Updates Q1 and Q2 projections accordingly.
\end{itemize}

\apimethod{Contains}{Check if a node exists}{bool Contains(string node)}
\item Returns \texttt{true} if the specified node is present in the tree.
\item Otherwise returns \texttt{false}.
\end{itemize}

\apimethod{GetParent}{Return the parent of a node}{string GetParent(string node)}
\item Returns the parent of \texttt{node}.
\item If \texttt{node} is the root, returns \texttt{null}.
\item Throws an error if the node is not found.
\end{itemize}

\section*{Structural Reconfiguration}

\apimethod{CutNPaste}{Detach a subtree from its root and attach it to another}{void CutNPaste(string u1, string u2)}
\item Extracts the subtree rooted at \texttt{u2} from the tree rooted at \texttt{u1}.
\item The subtree is detached from its original parent and reattached under a new one.
\item Ensures that no orphaned structure remains in the source tree.
\item Both source and destination trees must be structurally valid after the operation.
\end{itemize}

\apimethod{Delete}{Remove a node and reattach its children, if any, to its parent}{void Delete(string node)}
\item Deletes \texttt{node} from the tree.
\item If \texttt{node} has children, they are reattached to \texttt{RootOf(node)}.
\item Behaves as a shorthand for \texttt{CutNPaste(node, RootOf(node))}.
\item If \texttt{node} is the root, the operation is disallowed.
\end{itemize}

\section*{Accessors}

\apimethod{Q1}{Return the Q1 projection}{DList Q1}
\item Returns a copy of the Q1 projection (left-to-right traversal).
\item The projection reflects the current visible tree state.
\end{itemize}

\apimethod{Q2}{Return the Q2 projection}{DList Q2}
\item Returns a copy of the Q2 projection (right-to-left traversal).
\item Q2 is the sibling-reversed counterpart of Q1.
\end{itemize}

\apimethod{Source}{Return the symbolic source node}{string Source}
\item Returns the symbolic start marker of all projections.
\item This node is \textbf{not} part of the tree structure.
\item It conceptually acts as the left bracket delimiting Q1 and Q2.
\end{itemize}

\apimethod{Sink}{Return the symbolic sink node}{string Sink}
\item Returns the symbolic end marker of all projections.
\item This node is \textbf{not} part of the tree structure.
\item It conceptually acts as the right bracket delimiting Q1 and Q2.
\end{itemize}

\section*{Query Methods}

\subsection*{Reachability}

\apimethod{CanReach}{Check if one node can reach another}{bool CanReach(string u, string v)}
\item Returns \texttt{true} if there is a directed path from node \texttt{u} to node \texttt{v}.
\item The check is structural and respects tree hierarchy.
\end{itemize}

\apimethod{CanNotReach}{Check if one node cannot reach another}{bool CanNotReach(string u, string v)}
\item Returns \texttt{true} if there is no directed path from node \texttt{u} to node \texttt{v}.
\item Logical negation of \texttt{CanReach(u, v)}.
\end{itemize}

\subsection*{Ancestry and Structure}

\apimethod{IsRoot}{Check if a node is the root}{bool IsRoot(string node)}
\item Returns \texttt{true} if \texttt{node} is the root of the tree.
\end{itemize}

\apimethod{IsLeaf}{Check if a node has no children}{bool IsLeaf(string node)}
\item Returns \texttt{true} if \texttt{node} has no outgoing edges.
\end{itemize}

\apimethod{IsAncestor}{Check if one node is an ancestor of another}{bool IsAncestor(string u, string v)}
\item Returns \texttt{true} if node \texttt{u} lies on the path from the root to node \texttt{v}.
\end{itemize}

\apimethod{IsDescendant}{Check if one node is a descendant of another}{bool IsDescendant(string u, string v)}
\item Returns \texttt{true} if node \texttt{u} lies below node \texttt{v} in the tree.
\item Equivalent to \texttt{IsAncestor(v, u)}.
\end{itemize}

\apimethod{SameSubtree}{Check if two nodes belong to the same tree}{bool SameSubtree(string u, string v)}
\item Returns \texttt{true} if both nodes share the same root.
\item Useful for validating \texttt{CutNPaste} targets.
\end{itemize}

\apimethod{SubtreeSize}{Return the number of nodes in a subtree}{int SubtreeSize(string node)}
\item Returns the size (inclusive) of the subtree rooted at \texttt{node}.
\item Includes the node itself and all its descendants.
\end{itemize}

\apimethod{Depth}{Return the depth of a node}{int Depth(string node)}
\item Returns the distance from the root to \texttt{node}.
\item The root has depth 0.
\end{itemize}

\apimethod{PathToRoot}{Return the list of nodes from a node to the root}{List<string> PathToRoot(string node)}
\item Returns the sequence of nodes on the path from \texttt{node} to the root.
\item The last entry in the list is the root itself.
\end{itemize}

\apimethod{RootToPath}{Return the list of nodes from the root to a node}{List<string> RootToPath(string node)}
\item Returns the sequence of nodes on the path from the root to \texttt{node}.
\item The first entry in the list is the root itself.
\end{itemize}

\end{document}
